\section{One witness for a halting computation}\label{sec:computation}

Propositions~\ref{prop:power} and~\ref{prop:subset} provide single-fold
definitions of powers and binary containment. We now use these relations to
encode an entire computation by finitely many natural numbers. Determinism
will determine the computation, and a clock will determine the size of the
blocks used to encode it. Both features are needed: a fixed computation
admits many encodings if its block size is left free.
The method of recording register computations in a large power-of-two
base comes from Jones and Matiyasevich's bit-masking construction
\cite{jm1984}; see especially the exposition in
\cite[Section~4]{jm1991}. Determining the base from the input and
computation time already appears in Chaitin's account
\cite[Sections~2.3 and~2.6]{chaitin1987}. Here a clock and a dyadic
interval determine the base. We verify both the decoding of every
solution and the uniqueness of its complete auxiliary tuple.

\subsection{The machine model}

A counter machine has a finite set of states, a distinguished initial state
$q_*$, a distinguished halt state $h$, and finitely many counters with values
in $\NN$. Every state other than $h$ has exactly one instruction of one of
the following forms:
\begin{itemize}
\item $\operatorname{INC}(j)$: increase counter $j$ by one and go to a
specified state;
\item $\operatorname{TESTDEC}(j)$: if counter $j$ is zero, leave it zero and
go to a specified state; otherwise decrease it by one and go to a second
specified state.
\end{itemize}
The other counters remain unchanged. The halt state has no instruction.
Thus every configuration outside $h$ has exactly one successor. Time counts
individual instructions, including those used to implement a subroutine.
This is the classical counter instruction set of
Minsky~\cite[Theorem~Ia]{minsky1961}; register-machine formulations were
also developed by Shepherdson and Sturgis
\cite[Section~10, Result~10.3]{shepherdson1963}. We include the simulation needed
for the direct tuple inputs used here, so no encoded-input convention is
imported into the lemma.

\begin{lemma}\label{lem:counter-universality}
Let $n\geq 1$ and let $S\subseteq\NN^n$ be recursively enumerable. There
is a counter machine that, when its first $n$ counters are initialized to
$a_1,\ldots,a_n$ and its remaining counters to zero, reaches its halt state
if and only if $\mathbf a\in S$.
\end{lemma}

\begin{proof}
Choose a deterministic Turing machine that accepts precisely the unary
input words
$1^{a_1}\#\cdots1^{a_n}\#$ with $\mathbf a\in S$. Any rejecting terminal
behavior can be replaced by an infinite loop, so only acceptance halts.
We describe a simulation using the two instruction types above.

Store the tape strictly to the left and strictly to the right of the
current cell as two stacks; store the current symbol in the finite
control. Give the finitely many tape symbols distinct digits
$1,\ldots,q$, including a digit for the blank symbol, and put $D=q+1$.
A stack with successive symbols $e_0,e_1,\ldots,e_{m-1}$, starting at its
top, is stored as $\sum_{i<m}e_iD^i$. The empty stack is zero.
Allocate distinct counters for the two stacks, the input coordinates,
a scratch counter, and a counter that is always zero.

To push digit $e$ onto a stack of value $v$, compute $Dv+e$.
Repeatedly decrement the stack counter and, after each successful
decrement, increment the initially zero scratch counter exactly $D$
times. Transfer the scratch counter back by repeated decrements of it
and increments of the stack counter; then perform $e$ increments.
The scratch counter is again zero. All repetitions of a fixed length
are implemented by finitely many states.

To pop a nonempty stack, repeatedly decrement its counter, keeping the
number of successful decrements modulo $D$ in the finite control, and
increment the scratch counter after every $D$ successful decrements.
When the stack counter reaches zero, the scratch counter contains the
quotient by $D$ and the control contains the remainder. Transfer that
quotient back to the stack counter. Its original least significant
digit is nonzero, so the remainder gives the symbol popped. An empty
stack returns a blank and stays empty. Again the scratch counter ends
at zero. Moving the tape head pushes the written current symbol onto
one stack and pops the other. Storing additional blank symbols causes
no difficulty, since each finite computation uses finite stacks.

Initialize the right stack by pushing the input word in reverse: for
$i=n,n-1,\ldots,1$, push $\#$ and then push a $1$ for each successful
decrement of input counter $i$. Pop once to obtain the initial current
symbol; the left stack is empty. A zero test on the counter that is
always zero implements an unconditional jump. Every subroutine uses
finitely many states; duplicating these states supplies its finitely
many required return destinations. This gives a deterministic counter
machine with a sole halt state reached exactly on acceptance.
\end{proof}

Fix such a machine for the remainder of the construction. Write $r$ for
its number of counters and $s$ for its number of states. Its initial
counter values are $d_1,\ldots,d_r$, where each $d_j$ is an input
coordinate or zero. In the proof below it is harmless to allow any fixed
list of such initial values. Put
\begin{equation}\label{eq:machine-C}
 C=1+s+\sum_{j=1}^r d_j.
\end{equation}
The symbol $C$ denotes this expression, not an additional choice.
At time $t$ in any run, counter $j$ has value at most $d_j+t<C+t$,
since an instruction increases a counter by at most one. A clock with
value $C+t$ will therefore give a common bound for the counter values.

Regard an increment instruction as one directed arc, labeled $+j$.
Regard a test-decrement instruction as two directed arcs, labeled
$0j$ and $-j$. Each arc retains its specified destination. Let
$\mathcal A$ be this finite set of arcs, and let
$\mathcal A_j^+$ and $\mathcal A_j^-$ denote those labeled $+j$ and
$-j$, respectively. The halt state has no outgoing arc, and
$|\mathcal A|\leq 2s$.

\subsection{A finite system for the run}

\begin{proposition}\label{prop:machine-compiler}
Fix $n\ge1$ and a deterministic counter machine with finitely many states
and counters, instructions $\operatorname{INC}$ and $\operatorname{TESTDEC}$,
and one halt state. Initialize each counter by a fixed choice of an input
coordinate $a_i$ or zero. There is a finite system of integer polynomial
equations in $a\in\NN^n$ and natural auxiliary variables such that, for
every input $a$, the complete auxiliary tuple has exactly one solution
when the initialized machine halts, and no solution when it does not halt.
\end{proposition}

We first give the system using the single-fold relations $\Pow$ and
$\bits$. Every occurrence is a separate call with fresh auxiliary
variables, as in Section~\ref{sec:conventions}. Inequalities use the unique
slacks from that section. All sets indexing the displayed constraints
are fixed finite sets of states, arcs, or counters.

Introduce natural variables $w,B,b,k,H,J,M$ and require
\begin{align}
 &\Pow(2,w,B),\qquad B=4b,\qquad B>2C,\qquad k>0,
       \qquad\Pow(B,k,H),\label{eq:machine-base}\\
 &BJ=H,\qquad (B-1)M=H-1.\label{eq:machine-geometric}
\end{align}
These equations force
\begin{equation}\label{eq:machine-JM}
 B=2^w,\qquad J=B^{k-1},\qquad M=\sum_{t=0}^{k-1}B^t.
\end{equation}
The intended run has $k$ configurations, at times $0,\ldots,k-1$.
In its base-$B$ encodings, each time occupies one block of $w$ binary
digits. The number $M$ has a single $1$ in the lowest position of each
block.
Multiplication by $B$ shifts each block from time $t$ to time $t+1$.
The transition equations below use this shift to place a selected arc's
destination and its counter update in the next configuration.

For every state $q$, introduce $I_q$ and impose
\begin{equation}\label{eq:machine-state-partition}
 I_q\bits M,\qquad \sum_q I_q=M.
\end{equation}
The intended meaning is that the base-$B$ digit of $I_q$ at time $t$
is one exactly when the machine is in state $q$. For every test state
$q$, introduce $Z_q,N_q,V_q$ and impose
\begin{equation}\label{eq:machine-test-partition}
 Z_q\bits I_q,\qquad N_q\bits I_q,\qquad
 Z_q+N_q=I_q,\qquad Z_q+V_q=M.
\end{equation}
Here $Z_q$ records zero branches, $N_q$ records decrement branches, and
$V_q$ records the complement of the zero-branch times.
Define the arc expression $E_a$ to be $I_q$ for an increment arc from
$q$, and $Z_q$ or $N_q$ for the corresponding arc from a test state $q$.
These are expressions in existing variables. For every state impose
\begin{equation}\label{eq:machine-state-transition}
 I_q=\mathbf 1_{q=q_*}
       +B\sum_{\substack{a\in\mathcal A\\\operatorname{dest}(a)=q}}E_a.
\end{equation}
The coefficient $\mathbf 1_{q=q_*}$ is the fixed integer zero or one.

Counter digits will use only the lowest $w-1$ bits of each block,
so their values lie between zero and $B/2-1=2b-1$.
Reserving the top bit will let the decoding proof exclude signed
carries in the counter updates, which include subtraction.
For each counter $j=1,\ldots,r$, introduce $R_j,\rho_j,\eta_j$ and impose
\begin{align}
 &R_j\bits (2b-1)M,\qquad R_j=\rho_jJ+\eta_j,
       \qquad\eta_j<J,\label{eq:machine-counter-bound}\\
 &R_j=d_j+B\left(R_j-\rho_jJ+
       \sum_{a\in\mathcal A_j^+}E_a-
       \sum_{a\in\mathcal A_j^-}E_a\right).
       \label{eq:machine-counter-transition}
\end{align}
For every test state $q$ operating on counter $j$, additionally require
\begin{equation}\label{eq:machine-zero-test}
 R_j\bits (B-1)V_q.
\end{equation}
The intended base-$B$ digits of $R_j$ are the successive values of the
counter. The quotient $\rho_j$ is its final value. Removing $\rho_jJ$
before multiplying by $B$ makes the recurrence stop at the final time.
The mask $(B-1)V_q$ has a zero block at each time selected by $Z_q$
and has all $w$ binary digits equal to one in every other time block.
Thus the containment condition forces the tested counter to vanish at
the selected times.

Finally introduce $R_0,\rho_0,\eta_0$ for a clock and impose
\begin{align}
 &R_0\bits(2b-1)M,\qquad R_0=\rho_0J+\eta_0,
       \qquad\eta_0<J,\label{eq:machine-clock-bound}\\
 &R_0=C+B(R_0-\rho_0J+M-J),\label{eq:machine-clock-update}\\
 &b\leq\rho_0,\qquad 2\rho_0<B.\label{eq:machine-clock-scale}
\end{align}
Its intended value at time $t$ is $C+t$. The last two inequalities put
its final value in $[B/4,B/2)$; this leaves room for all counter digits
and will select exactly one power-of-two base. We now prove that every
solution has these intended meanings, and that every halting run supplies
a solution.

\begin{proof}[Proof of Proposition~\ref{prop:machine-compiler}]
\emph{States and the final halt.}
Take any solution. Since $B>2C$ and $C\geq 1+s$, we have $B>2s$;
also $B=4b\geq 4$. By \eqref{eq:machine-JM}, binary containment in $M$
means that every base-$B$ digit of $I_q$ is zero or one, with no
nonzero digits outside times $0,\ldots,k-1$. The sum in
\eqref{eq:machine-state-partition} has at most $s<B$ in each digit,
so it has no carries. Thus exactly one state is present at each time.
Similarly, \eqref{eq:machine-test-partition} partitions the times of
each test state into its two branches. Consequently a nonhalt state
has exactly one selected outgoing arc, and the halt state has none.

The digit sums on the right of
\eqref{eq:machine-state-transition} are at most
$|\mathcal A|<B$. Hence those equations also have no carries. They
give state $q_*$ at time zero, and at every later time they give the
destination of the selected arc at the preceding time. The halt mask
has an exact form. Each nonhalt state occurrence selects one arc, so
$\sum_{a\in\mathcal A}E_a=M-I_h$. Summing the state equations gives
\[
 M=1+B(M-I_h),\qquad
 BI_h=1+(B-1)M=H=BJ.
\]
Thus $I_h=J$: the halt state occurs precisely at time $k-1$.
There is exactly one selected arc at each time $t<k-1$, and none
at the final time. This also excludes earlier halts and padded records.

\emph{Counter values and zero tests.}
Since $B=2^w=4b$, the integer
$(2b-1)M=(B/2-1)M$ allows exactly the lowest $w-1$ bits of each
block. Thus we can write uniquely
\begin{equation}\label{eq:machine-counter-digits}
 R_j=\sum_{t=0}^{k-1}r_{j,t}B^t,
 \qquad 0\leq r_{j,t}\leq B/2-1.
\end{equation}
Division by $J=B^{k-1}$ gives $\rho_j=r_{j,k-1}$ and determines
$\eta_j$ uniquely. Let $a_{j,t}$ and $f_{j,t}$ indicate whether the
selected arc at time $t$ increments or decrements counter $j$.
They belong to $\{0,1\}$, never both equal one, and are both zero
at the final time.

Expand \eqref{eq:machine-counter-transition} using these digits.
The coefficient of $B^0$ in the resulting signed sum is
$r_{j,0}-d_j$. For $1\leq t<k$ it is
\begin{equation}\label{eq:machine-signed-coefficient}
 r_{j,t}-r_{j,t-1}-a_{j,t-1}+f_{j,t-1}.
\end{equation}
There are no higher coefficients: the final digit of $R_j$ was
removed, and no final arc was selected. The constant coefficient has
absolute value less than $B$, because $d_j<C<B/2$; every coefficient
in \eqref{eq:machine-signed-coefficient} has absolute value at most
$B/2<B$. An integer sum $\sum_{t=0}^{k-1}c_tB^t=0$ with
$|c_t|<B$ has every $c_t=0$: reduce modulo $B$, use that the only
multiple of $B$ in $(-B,B)$ is zero, divide by $B$, and repeat.
It follows that
\begin{equation}\label{eq:machine-actual-update}
 r_{j,0}=d_j,\qquad
 r_{j,t+1}=r_{j,t}+a_{j,t}-f_{j,t}\quad(0\leq t<k-1).
\end{equation}
This argument rules out borrowing as well as carrying.

Since $Z_q$ selects some of the ones of $M$, the equation
$V_q=M-Z_q$ gives base-$B$ digit zero at its zero-branch times and
one at all other times. Thus $(B-1)V_q$ allows every bit in a block
except at those zero-branch times, where it allows none.
Condition~\eqref{eq:machine-zero-test} forces the tested counter
to be zero whenever the zero branch is selected. On a decrement
branch, \eqref{eq:machine-actual-update} and nonnegativity of the
next digit force the tested counter to be positive. Hence every
test and every update is a valid machine instruction.

\emph{The clock and the block size.}
Apply the same coefficient argument to the clock. In
\eqref{eq:machine-clock-update}, the expression $M-J$ has digit
one exactly at times $t<k-1$. Its digit recurrence is therefore
$r_{0,0}=C$ and $r_{0,t+1}=r_{0,t}+1$. In particular,
\begin{equation}\label{eq:machine-clock-value}
 \rho_0=C+k-1.
\end{equation}
The state and counter arguments show that the encoded configurations
form the actual deterministic run from the specified initialization,
and that it first reaches $h$ at time $k-1$. If the machine does not
halt, no solution exists. If it halts after $T$ instructions, then
$k=T+1$ is forced. Conditions~\eqref{eq:machine-clock-scale} become
\begin{equation}\label{eq:machine-unique-scale}
 2(C+T)<B\leq4(C+T).
\end{equation}
For every positive integer $L$, the interval $(2L,4L]$ contains
exactly one power of two: take the smallest power strictly greater
than $2L$, whose predecessor is at most $2L$. Its successor exceeds
$4L$. Thus \eqref{eq:machine-unique-scale} determines $B$ uniquely,
including when an endpoint is a power of two.

\emph{Existence from a halting run.}
Suppose the initialized machine halts after $T$ instructions. Set
$k=T+1$ and choose the unique power of two $B$ in
\eqref{eq:machine-unique-scale}. Since $C\geq2$, we have $B>4$,
so $B$ is divisible by four. Set $b=B/4$, choose the unique exponent
$w$ with $B=2^w$, set $H=B^k$, and use
\eqref{eq:machine-JM} to define $J,M$.
At time $t\leq T$ each counter has value at most $d_j+t$, because
one instruction increases a counter by at most one. Hence its value
is at most $\sum_jd_j+T<C+T<B/2$. The clock values $C+t$ also
are less than $B/2$. Form $R_j$ and $R_0$ from these digits, and
form the state and branch masks from the actual run.

All state equations follow from the transitions. The zero-test
constraints follow from the zero branches. Summing the actual
counter updates gives \eqref{eq:machine-counter-transition},
with the final digit removed before shifting; summing the clock
updates gives \eqref{eq:machine-clock-update}. The remaining
inequalities follow from \eqref{eq:machine-unique-scale}.
Every power and binary-containment call is true on these arguments
and therefore has its required unique auxiliary tuple.
This also handles $T=0$: then $q_*=h$, $k=1$, $J=M=1$,
all arc masks vanish, and every update sum is empty.

\emph{Uniqueness of the complete tuple.}
In any solution, the deterministic run fixes $T$, every state,
every selected branch, and every counter value. It therefore fixes
$k$; the clock inequalities then fix $B$, and hence $w,b,H,J,M$.
The run fixes all $I_q,Z_q,N_q,R_j,R_0$, while
$V_q=M-Z_q$ fixes every $V_q$. Division by the positive $J$
fixes every $\rho_j,\eta_j$, including the clock pair.
These are all the displayed variables. Each introduced argument
variable is equal to an already determined polynomial expression;
every inequality slack is determined by its endpoints. Finally,
Propositions~\ref{prop:power} and~\ref{prop:subset} determine
all auxiliary variables in every expanded call. Their internal
case selectors and inactive variables are included in that
guarantee. Thus uniqueness holds for the complete tuple, not only
for the numbers recording the run.
\end{proof}

\subsection{The resulting polynomial}

\begin{proof}[Proof of Theorem~\ref{thm:main}]
Given the recursively enumerable set $S\subseteq\NN^n$, choose the
machine of Lemma~\ref{lem:counter-universality} and apply
Proposition~\ref{prop:machine-compiler}. Expand each call to a
single-fold relation into its finite polynomial system, using fresh
auxiliary variables, and expand all inequalities and disjoint cases
according to Section~\ref{sec:conventions}.
The number of equations and variables is finite: the machine is
fixed, its state and counter sets are finite, and each relation call
has a fixed finite definition. The time parameter $k$ is one natural
variable; the digit expansions used in the proof introduce no
additional variables indexed by its value.

Write the resulting integer polynomial equations as
$F_1(\mathbf a,\mathbf z)=\cdots=F_N(\mathbf a,\mathbf z)=0$,
where $\mathbf z$ includes every auxiliary coordinate, and set
\[
 P(\mathbf a,\mathbf z)=\sum_{i=1}^N F_i(\mathbf a,\mathbf z)^2.
\]
On natural tuples, $P=0$ is equivalent to the simultaneous
vanishing of all $F_i$. This operation introduces no variables and
preserves the solution set exactly. By the compiler proposition,
that set has one element when $\mathbf a\in S$ and is empty
otherwise, as asserted.
\end{proof}
